\documentclass[
	% -- opções da classe memoir --
	12pt,				% tamanho da fonte
	openright,			% capítulos começam em pág ímpar (insere página vazia caso preciso)
	twoside,			% para impressão em recto e verso. Oposto a oneside
	a4paper,			% tamanho do papel. 
	% -- opções da classe abntex2 --
	%chapter=TITLE,		% títulos de capítulos convertidos em letras maiúsculas
	%section=TITLE,		% títulos de seções convertidos em letras maiúsculas
	%subsection=TITLE,	% títulos de subseções convertidos em letras maiúsculas
	%subsubsection=TITLE,% títulos de subsubseções convertidos em letras maiúsculas
	% -- opções do pacote babel --
	english,			% idioma adicional para hifenização
	brazilian			% o último idioma é o principal do documento
	]{abntex2}

% ---
% Pacotes básicos 
% ---
\usepackage{lmodern}			% Usa a fonte Latin Modern			
\usepackage[T1]{fontenc}		% Selecao de codigos de fonte.
\usepackage[utf8]{inputenc}		% Codificacao do documento (conversão automática dos acentos)
\usepackage{indentfirst}		% Indenta o primeiro parágrafo de cada seção.
\usepackage{color}				% Controle das cores
\usepackage[table]{xcolor}		% Cores em tabelas
\usepackage{graphicx}			% Inclusão de gráficos
\usepackage{microtype} 			% para melhorias de justificação
\usepackage{amsmath, amssymb}	% Pacotes para cálculos matemáticos e símbolos
\usepackage{array}				% Melhorias em tabelas
\usepackage{xurl}				% Quebra de linha automática em URLs longas
\raggedbottom					% Evita Underfull \vbox em páginas parcialmente preenchidas
% ---

% ---
% Pacotes de citações
% ---
\usepackage[brazilian,hyperpageref]{backref}	 % Paginas com as citações na bibl
\usepackage[alf]{abntex2cite}	% Citações padrão ABNT

% --- 
% CONFIGURAÇÕES DE PACOTES
% --- 

% ---
% Configurações do pacote backref
% Usado sem a opção hyperpageref de backref
\renewcommand{\backrefpagesname}{Citado na(s) página(s):~}
% Texto padrão antes do número das páginas
\renewcommand{\backref}{}
% Define os textos da citação
\renewcommand*{\backrefalt}[4]{
	\ifcase #1 %
		Nenhuma citação no texto.%
	\or
		Citado na página #2.%
	\else
		Citado #1 vezes nas páginas #2.%
	\fi}%
% ---

% ---
% Informações de dados para CAPA e FOLHA DE ROSTO
% ---
\titulo{Análise e Comparação de Estratégias Algorítmicas}
\autor{Integrante 1 \\ Integrante 2 \\ Integrante 3 \\ Integrante 4 \\ Integrante 5 \\ Integrante 6 \\ Integrante 7}
\local{Juazeiro -- BA}
\data{2026}
\orientador{Prof. Alexandre Freire de Souza}
\instituicao{%
  Universidade Federal do Vale do São Francisco -- UNIVASF
  \par
  Colegiado de Engenharia da Computação / Ciência da Computação}
\tipotrabalho{Relatório Técnico}
% O preambulo deve conter o tipo do trabalho, o objetivo, 
% o nome da instituição e a área de concentração 
\preambulo{Relatório técnico apresentado à disciplina de Análise de Algoritmos da Universidade Federal do Vale do São Francisco (UNIVASF) como requisito avaliativo para estudo, análise e comparação de estratégias algorítmicas.}
% ---

% ---
% Configurações de aparência do PDF final

% alterando o aspecto da cor azul
\definecolor{blue}{RGB}{41,5,195}

% informações do PDF
\makeatletter
\hypersetup{
     	%pagebackref=true,
		pdftitle={\@title}, 
		pdfauthor={Equipe},
    	pdfsubject={\imprimirpreambulo},
	    pdfcreator={LaTeX with abnTeX2},
		pdfkeywords={análise de algoritmos}{heurísticas}{busca exaustiva}{algoritmos de aproximação}{TSP}, 
		colorlinks=true,       		% false: boxed links; true: colored links
    	linkcolor=blue,          	% color of internal links
    	citecolor=blue,        		% color of links to bibliography
    	filecolor=magenta,      		% color of file links
		urlcolor=blue,
		bookmarksdepth=4
}
\makeatother
% --- 

% ---
% Posiciona figuras e tabelas no topo da página quando adicionadas sozinhas
% em um página em branco. Ver https://github.com/abntex/abntex2/issues/170
\makeatletter
\setlength{\@fptop}{5pt} % Set distance from top of page to first float
\makeatother
% ---

% --- 
% Espaçamentos entre linhas e parágrafos 
% --- 

% O tamanho do parágrafo é dado por:
\setlength{\parindent}{1.3cm}

% Controle do espaçamento entre um parágrafo e outro:
\setlength{\parskip}{0.2cm}  % tente também \onelineskip

% ----
% Início do documento
% ----
\begin{document}

% Seleciona o idioma do documento (conforme pacotes do babel)
\selectlanguage{brazilian}

% Retira espaço extra obsoleto entre as frases.
\frenchspacing 

% ----------------------------------------------------------
% ELEMENTOS PRÉ-TEXTUAIS
% ----------------------------------------------------------
% \pretextual

% ---
% Capa
% ---
\imprimircapa
% ---

% ---
% Folha de rosto
% ---
\imprimirfolhaderosto{}
% ---

% ---
% inserir o sumario
% ---
\pdfbookmark[0]{\contentsname}{toc}
\tableofcontents*
\cleardoublepage
% ---

% ----------------------------------------------------------
% ELEMENTOS TEXTUAIS
% ----------------------------------------------------------
\textual

% ----------------------------------------------------------
% Introdução (capítulo sem numeração, mas presente no Sumário)
% ----------------------------------------------------------
\chapter*[Introdução]{Introdução}
\addcontentsline{toc}{chapter}{Introdução}
% ----------------------------------------------------------

\section*{Objetivo Geral e Assuntos Abordados}

\section*{Justificativa: Tempo de Processamento e Crescimento Assintótico}

Na Ciência da Computação, a dificuldade de resolver um problema à medida que o volume de dados cresce é explicada pela \textbf{Teoria da Complexidade Computacional} e pela \textbf{Análise Assintótica de Algoritmos}. O fator determinante não é apenas o poder de processamento do \textit{hardware}, mas sim a \textbf{taxa de crescimento} da função matemática $T(n)$ que relaciona o tamanho da entrada ($n$) à quantidade de operações elementares que o algoritmo precisa executar no pior caso.

\subsection*{Algoritmos Polinomiais versus Exponenciais e Fatoriais}

Para compreender por que determinados problemas se tornam intratáveis, divide-se o comportamento assintótico dos algoritmos em duas grandes categorias:

\begin{itemize}
    \item \textbf{Crescimento Polinomial (Problemas Tratáveis -- Classe P):} Ocorre quando o tempo de execução é limitado superiormente por um polinômio da forma $O(n^k)$, onde $k$ é uma constante positiva (ex.: busca linear $O(n)$, ordenação \textit{MergeSort} $O(n \log n)$ ou matrizes em $O(n^2)$). Nesses casos, o aumento de dados provoca um crescimento previsível e gerenciável no tempo de processamento. Por exemplo, em um algoritmo linear $O(n)$, dobrar o volume de dados apenas dobra o tempo de espera; em um quadrático $O(n^2)$, dobrar a entrada quadruplica o tempo.
    \item \textbf{Crescimento Superpolinomial, Exponencial ou Fatorial (Problemas Intratáveis em Larga Escala):} Ocorre quando a variável $n$ encontra-se no expoente, como em $O(2^n)$, ou envolve permutações completas, como em $O(n!)$. Nesses problemas (frequentemente pertencentes à classe \textbf{NP-difícil}), a adição de um único novo elemento aos dados de entrada não soma um valor fixo ao tempo de execução, mas sim \textbf{multiplica} todo o esforço computacional anterior.
\end{itemize}

\subsection*{A Relação Direta com o Tempo de Processamento}

Para um tamanho de entrada muito pequeno (por exemplo, $n = 5$), a diferença prática de tempo entre um algoritmo eficiente de complexidade $O(n^2)$ e um algoritmo exaustivo de complexidade $O(2^n)$ é imperceptível para um ser humano, pois ambos executam em frações de microssegundo ($25$ operações contra $32$ operações). 

Contudo, à medida que $n$ cresce, ocorre uma divergência drástica: o número de estados ou combinações possíveis no espaço de busca aumenta mais rápido do que qualquer capacidade física de processamento. Considere um computador hipotético capaz de realizar \textbf{1 bilhão ($10^9$) de operações por segundo}. A \autoref{tab:crescimento_assintotico} ilustra como o tempo necessário para encontrar uma solução evolui conforme a quantidade de dados ($n$) aumenta em diferentes ordens de complexidade:

\begin{table}[htb]
\centering
\caption{\label{tab:crescimento_assintotico}Relação entre o tamanho da entrada ($n$) e o tempo de execução para um processador de $10^9$ operações/segundo.}
\renewcommand{\arraystretch}{1.3}
\begin{tabular}{|l|c|c|c|c|}
\hline
\rowcolor{gray!20}
\textbf{Complexidade} & \textbf{$n = 10$} & \textbf{$n = 20$} & \textbf{$n = 50$} & \textbf{$n = 100$} \\ \hline
$O(n \log_2 n)$ & $0,03\ \mu\text{s}$ & $0,08\ \mu\text{s}$ & $0,28\ \mu\text{s}$ & $0,66\ \mu\text{s}$ \\ \hline
$O(n^2)$ & $0,1\ \mu\text{s}$ & $0,4\ \mu\text{s}$ & $2,5\ \mu\text{s}$ & $10\ \mu\text{s}$ \\ \hline
$O(n^3)$ & $1\ \mu\text{s}$ & $8\ \mu\text{s}$ & $125\ \mu\text{s}$ & $1\ \text{ms}$ \\ \hline
$O(2^n)$ & $1,02\ \mu\text{s}$ & $1,05\ \text{ms}$ & $13\ \text{dias}$ & $4 \times 10^{13}\ \text{anos}$ \\ \hline
$O(n!)$ & $3,6\ \text{ms}$ & $77,1\ \text{anos}$ & $> 10^{48}\ \text{anos}$ & Inviável \\ \hline
\end{tabular}
\end{table}

Observa-se que, enquanto um algoritmo cúbico $O(n^3)$ resolve uma instância de $n = 100$ elementos em apenas 1 milissegundo, um algoritmo de complexidade exponencial $O(2^n)$ salta de cerca de 1 milissegundo (para $n = 20$) para 13 dias (para $n = 50$) e atinge dezenas de trilhões de anos para $n = 100$ (superando em milhares de vezes a idade estimada do universo). 

Portanto, a dificuldade computacional não decorre de uma limitação tecnológica momentânea dos computadores atuais, mas de uma barreira matemática: em problemas cuja solução exata exige explorar espaços de busca exponenciais ou fatoriais, mesmo que a velocidade dos processadores fosse multiplicada por mil, o ganho prático permitiria aumentar o tamanho da entrada em apenas algumas unidades antes de esbarrar novamente na barreira do tempo impraticável.

% ----------------------------------------------------------
% PARTE I
% ----------------------------------------------------------
\part{Estratégias de Resolução e Busca}
% ----------------------------------------------------------

% ----------------------------------------------------------
% Capítulo 1 (Integrante 1)
% ----------------------------------------------------------
\chapter{Problema do Caixeiro-Viajante e Algoritmos Gulosos}

\section{Conceito, Definição e Representação do Problema do Caixeiro-Viajante}

O \textbf{Problema do Caixeiro-Viajante} (\textit{Traveling Salesman Problem} -- TSP) é um dos problemas mais clássicos da Ciência da Computação e da Otimização Combinatória.

\subsection{Representação por Meio de Grafos e Distâncias entre Cidades}

Para que um sistema computacional possa processar e resolver esse problema, é necessário modelá-lo matematicamente por meio da Teoria dos Grafos, estruturando um grafo ponderado $G = (V, E)$. As informações essenciais incluem:

\begin{itemize}
    \item \textbf{Conjunto de localidades (Vértices $V$):} A lista exata das cidades que devem ser visitadas, bem como a definição clara de qual delas atua como ponto de partida (e de retorno, caso o veículo precise voltar ao centro de distribuição de origem).
    \item \textbf{Conexões viárias (Arestas $E$):} O mapeamento das estradas ou rotas diretas que interligam cada par de cidades, indicando se o grafo é completo (todas as cidades se conectam entre si) ou se existem restrições de acesso.
    \item \textbf{Matriz de distâncias ou custos (Pesos $w$):} O valor numérico correspondente à distância física (em quilômetros) entre cada par de cidades $(u, v)$. Em cenários logísticos reais, essa matriz pode ser assimétrica (grafo direcionado), visto que a distância de $A$ para $B$ pode divergir da distância de $B$ para $A$ em razão de vias de mão única, desvios ou topografia.
    \item \textbf{Restrições operacionais adicionais:} Dados de regras de negócio, tais como a obrigatoriedade de visitar cada cidade exatamente uma única vez (formando um ciclo hamiltoniano) e eventuais bloqueios de tráfego ou limites de peso nas rodovias.
\end{itemize}

\subsection{A Dificuldade Computacional de Encontrar a Rota Ótima}

Encontrar a melhor rota garantida (solução ótima global) é uma tarefa computacionalmente árdua devido ao fenômeno conhecido como \textbf{explosão combinatória}. Na teoria da complexidade computacional, a versão de otimização do Problema do Caixeiro-Viajante pertence à classe dos problemas \textbf{NP-difíceis} (\textit{NP-hard}), o que significa que não se conhece nenhum algoritmo capaz de resolvê-lo de forma exata em tempo polinomial para qualquer instância.

Se um computador tentar solucionar o problema por força bruta (busca exaustiva), avaliando todas as permutações possíveis para identificar a menor distância, o número de rotas distintas para $n$ cidades (fixando-se a cidade inicial e assumindo distâncias simétricas entre ida e volta) é dado pela expressão matemática:

\begin{equation}
    R = \frac{(n - 1)!}{2}
\end{equation}

Aplicando essa fórmula para um cenário ilustrativo de $n = 20$ cidades de uma empresa de entregas, temos:

\begin{equation}
    R = \frac{(20 - 1)!}{2} = \frac{19!}{2} = \frac{121.645.100.408.832.000}{2} = 60.822.550.204.416.000\text{ rotas}
\end{equation}

Existem mais de \textbf{60,8 quatrilhões} de rotas possíveis (ou cerca de $1,21 \times 10^{17}$ caso as distâncias de ida e volta sejam diferentes). Mesmo que um processador moderno conseguisse avaliar 1 bilhão ($10^9$) de rotas completas por segundo, ele levaria cerca de $60.822.550$ segundos --- o que equivale a aproximadamente \textbf{704 dias} (quase dois anos de processamento ininterrupto) --- apenas para testar todas as possibilidades de uma única entrega de 20 cidades. Embora técnicas avançadas de programação dinâmica (como o algoritmo de Held-Karp, de complexidade $O(n^2 \cdot 2^n)$) consigam resolver 20 cidades em frações de segundo, basta aumentar a escala logística para 60 ou 100 pontos de entrega para que até os algoritmos exatos mais sofisticados se tornem inviáveis em tempo hábil.

\section{Introdução às Estratégias Heurísticas}

Diante da barreira de processamento imposta pelos problemas de natureza exponencial e fatorial (como os pertencentes à classe NP-difícil), a Ciência da Computação e a Pesquisa Operacional recorrem a métodos alternativos aos algoritmos exatos. Entre as abordagens mais difundidas na prática estão os métodos heurísticos.

\subsection{Definição de Heurística na Computação}

Etimologicamente, o termo \textit{heurística} deriva do verbo grego \textit{heuriskein}, que significa ``descobrir'' ou ``encontrar''. Na resolução de problemas computacionais, uma heurística é uma técnica, estratégia ou conjunto de regras práticas (\textit{rules of thumb}) projetadas para guiar a busca por uma solução viável dentro de um espaço de estados muito vasto, utilizando conhecimento específico do domínio do problema ou critérios intuitivos de decisão.

Ao contrário de um algoritmo exato (que explora sistematicamente o espaço de busca até comprovar matematicamente a otimalidade do resultado), a heurística atua como um ``atalho cognitivo'' ou computacional: ela aplica critérios simplificadores (como escolhas gulosas, podas agressivas na árvore de decisão ou buscas locais por melhoria incremental) para descartar a imensa maioria das combinações possíveis e convergir diretamente para uma resposta promissora.

\subsection{Principal Finalidade de Utilização}

A finalidade primordial de empregar uma heurística é \textbf{viabilizar a obtenção de soluções de boa qualidade em um tempo computacional aceitável (polinomial)} e com baixo consumo de memória, especialmente para problemas complexos ou de larga escala nos quais os métodos exatos falham por exaustão de tempo.

Em termos de projeto de algoritmos, o uso de uma heurística representa uma troca deliberada (\textit{trade-off}) entre \textbf{otimalidade} e \textbf{eficiência (velocidade)}. Quando uma aplicação real exige decisões em tempo real --- como o cálculo de uma rota em um aplicativo de GPS, o escalonamento de pacotes em um roteador de rede ou a movimentação de um agente em um jogo eletrônico ---, a heurística permite entregar uma resposta funcional em milissegundos, contornando a explosão combinatória.

\subsection{Garantia de Otimalidade da Solução}

\textbf{Uma heurística não garante que a solução encontrada será a melhor solução possível (solução ótima global)}. Além disso, em sua definição estrita, uma heurística pura também não oferece uma garantia matemática formal sobre o quão próxima ou distante a solução obtida está do valor ótimo verdadeiro.

Isso ocorre porque, para ganhar velocidade, o método heurístico toma decisões baseadas em informações parciais ou locais (visão míope do problema) e ignora regiões inteiras do espaço de busca. Esse comportamento sujeita o algoritmo a dois fenômenos clássicos:

\begin{enumerate}
    \item \textbf{Estagnação em Ótimos Locais:} Ao buscar melhorias passo a passo, a heurística pode atingir uma solução que é a melhor dentro de sua vizinhança imediata (um \textit{ótimo local}), mas que é significativamente inferior à melhor solução existente em outra região inexplorada do espaço de estados (o \textit{ótimo global}).
    \item \textbf{Propagação de Escolhas Iniciais Subótimas:} Em heurísticas construtivas (que montam a solução elemento por elemento), uma decisão que parece altamente vantajosa nas primeiras etapas pode restringir severamente as opções futuras, forçando o algoritmo a aceitar escolhas muito ruins no final do processo.
\end{enumerate}

\section{Funcionamento de um Algoritmo Guloso Aplicado ao TSP}

Uma das abordagens mais conhecidas na literatura de Otimização Combinatória e Análise de Algoritmos para o TSP é a \textbf{Heurística do Vizinho Mais Próximo} (\textit{Nearest Neighbor Heuristic} -- NNH). Trata-se de uma heurística construtiva clássica baseada no paradigma de \textbf{algoritmos gulosos} (\textit{greedy algorithms}), amplamente empregada para gerar soluções rápidas para o Problema do Caixeiro-Viajante.

\subsection{Etapas da Estratégia}

O algoritmo constrói a rota passo a passo, mantendo um conjunto de vértices (cidades) já visitados e outro de vértices pendentes. O seu funcionamento divide-se em quatro etapas sequenciais:

\begin{enumerate}
    \item \textbf{Inicialização:} Define-se a cidade de origem $v_0$ como a localização atual do entregador e marca-se $v_0$ como visitada.
    \item \textbf{Seleção Gulosa:} A partir da cidade atual $u$, o algoritmo inspeciona todas as arestas que conectam $u$ às cidades ainda não visitadas e seleciona a cidade $v$ cujo custo de deslocamento $d(u, v)$ seja o menor possível.
    \item \textbf{Atualização:} O entregador desloca-se para a cidade $v$, que passa a ser a nova localização atual e é marcada como visitada.
    \item \textbf{Critério de Parada e Fechamento:} Os passos 2 e 3 são repetidos iterativamente até que todas as $n$ cidades tenham sido visitadas. Por fim, adiciona-se uma última aresta conectando a última cidade visitada de volta ao ponto de partida $v_0$, fechando o ciclo hamiltoniano.
\end{enumerate}

\subsection{Eficiência Computacional e Limitações da Escolha Gulosa}

A principal vantagem dessa abordagem é a sua \textbf{extrema eficiência computacional e simplicidade de implementação}. Enquanto a busca exaustiva pela rota perfeita exige avaliar $(n-1)!/2$ permutações (complexidade fatorial $O(n!)$), a Heurística do Vizinho Mais Próximo reduz drasticamente o esforço para um \textbf{tempo polinomial quadrático}, de ordem $O(n^2)$.

Isso ocorre porque, na primeira cidade, o algoritmo faz $n-1$ comparações de distância; na segunda, faz $n-2$ comparações, e assim sucessivamente, totalizando:

\begin{equation}
    C(n) = (n - 1) + (n - 2) + \dots + 1 = \frac{n(n - 1)}{2} \in O(n^2)
\end{equation}

Na prática, para uma rota com $1.000$ cidades, um computador executa cerca de $500.000$ comparações simples e entrega uma rota completa em poucos milissegundos, sem necessidade de estruturas complexas de memória ou de \textit{backtracking} (retrocesso).

Entretanto, \textbf{essa estratégia não garante que a distância total percorrida será a menor possível}. A justificativa reside na \textbf{natureza míope (localmente ótima)} do critério guloso: em cada etapa, o entregador toma uma decisão olhando exclusivamente para o ganho imediato (a menor aresta disponível naquele instante), ignorando completamente a topologia global do mapa e as consequências futuras dessa escolha. Ao não prever o que acontecerá no restante do percurso, o algoritmo frequentemente deixa cidades geograficamente isoladas (ou em extremos opostos do mapa) ``esquecidas'' para o final. Quando restam poucas opções, o entregador torna-se obrigado a percorrer arestas extremamente longas para visitar essas cidades remanescentes e retornar à base, anulando toda a economia obtida nos passos iniciais.

\section{Exemplo Prático com Quatro Cidades}

Para ilustrar o funcionamento passo a passo e as limitações dessa estratégia, considere um entregador que parte da matriz na cidade \textbf{A} e precisa atender às cidades \textbf{B}, \textbf{C} e \textbf{D}, retornando a \textbf{A} no final. As distâncias em quilômetros entre as quatro localidades estão dispostas na matriz simétrica apresentada na \autoref{tab:matriz_tsp_4cidades}:

\begin{table}[htb]
\centering
\caption{\label{tab:matriz_tsp_4cidades}Matriz de distâncias (em km) para o cenário de quatro cidades.}
\renewcommand{\arraystretch}{1.2}
\begin{tabular}{|c|c|c|c|c|}
\hline
\rowcolor{gray!20}
\textbf{De / Para} & \textbf{Cidade A} & \textbf{Cidade B} & \textbf{Cidade C} & \textbf{Cidade D} \\ \hline
\textbf{Cidade A} & $0$ & $10$ & $12$ & $60$ \\ \hline
\textbf{Cidade B} & $10$ & $0$ & $15$ & $12$ \\ \hline
\textbf{Cidade C} & $12$ & $15$ & $0$ & $100$ \\ \hline
\textbf{Cidade D} & $60$ & $12$ & $100$ & $0$ \\ \hline
\end{tabular}
\end{table}

Analisando o comportamento da \textbf{Heurística do Vizinho Mais Próximo} em comparação com a \textbf{Rota Ótima}:

\begin{itemize}
    \item \textbf{Caminho escolhido pela Heurística ($A \rightarrow B \rightarrow D \rightarrow C \rightarrow A$):}
    \begin{enumerate}
        \item Partindo de \textbf{A}, as distâncias são: $B=10$, $C=12$ e $D=60$. O entregador escolhe a mais próxima: \textbf{B} ($10\text{ km}$).
        \item Estando em \textbf{B}, as não visitadas são $C=15$ e $D=12$. Ele escolhe a mais próxima: \textbf{D} ($12\text{ km}$).
        \item Estando em \textbf{D}, a única cidade não visitada restante é \textbf{C}. Ele é forçado a ir de \textbf{D} para \textbf{C}, percorrendo uma estrada muito longa de $100\text{ km}$ (a ``armadilha'' deixada para o final).
        \item De \textbf{C}, ele retorna à origem \textbf{A} ($12\text{ km}$).
        \item \textbf{Distância Total da Heurística:} $10 + 12 + 100 + 12 = \mathbf{134\text{ km}}$.
    \end{enumerate}
    \item \textbf{Caminho da Rota Ótima ($A \rightarrow C \rightarrow B \rightarrow D \rightarrow A$):}
    \begin{enumerate}
        \item Se no início o entregador abdicasse do vizinho imediatamente mais próximo ($B=10$) e fosse primeiro para \textbf{C} ($12\text{ km}$), de lá seguiria para \textbf{B} ($15\text{ km}$), depois para \textbf{D} ($12\text{ km}$) e retornaria direto a \textbf{A} ($60\text{ km}$), evitando a aresta crítica de $100\text{ km}$ entre C e D.
        \item \textbf{Distância Total Ótima:} $12 + 15 + 12 + 60 = \mathbf{99\text{ km}}$.
    \end{enumerate}
\end{itemize}

Nesse cenário concreto, a sedução de economizar apenas $2\text{ km}$ no primeiro movimento ($A \rightarrow B$ em vez de $A \rightarrow C$) penalizou o entregador no final da rota, gerando um percurso total $35\text{ km}$ mais longo ($35,3\%$ pior que o ideal). O mesmo fenômeno ocorre quando cidades estão distribuídas ao longo de uma linha reta com distâncias ligeiramente crescentes em lados alternados: o entregador fica fazendo um ``zigue-zague'' cada vez mais amplo de um extremo ao outro, em vez de percorrer um lado inteiro primeiro e o outro na volta.

% ----------------------------------------------------------
% Capítulo 2 (Integrante 2)
% ----------------------------------------------------------
\chapter{Busca Local e Estudo de Caso no TSP}

\section{Conceito de Busca Local}

\section{Funcionamento de Estratégias de Melhoria de Soluções}

\section{Exemplos de Operações de Troca ou Reorganização de Rotas}

\section{Aplicação de uma Estratégia de Busca Local ao Exemplo do TSP}

\section{Comparação entre a Solução Inicial e a Solução Após a Tentativa de Melhoria}

\section{Limitações das Estratégias de Busca Local}

% ----------------------------------------------------------
% Capítulo 3 (Integrante 3)
% ----------------------------------------------------------
\chapter{Algoritmos Exponenciais: Backtracking}

\section{Conceito de Backtracking}

\section{Funcionamento da Exploração do Espaço de Busca}

\section{Construção de uma Árvore de Possibilidades}

\section{Aplicação do Backtracking em um Problema Computacional}

\section{Apresentação de Pseudocódigo Simplificado}

\section{Discussão sobre o Custo Computacional e Crescimento Exponencial}

% ----------------------------------------------------------
% Capítulo 4 (Integrante 4)
% ----------------------------------------------------------
\chapter[Algoritmos Exponenciais: Branch-and-Bound]{Algoritmos Exponenciais:\\ Branch-and-Bound}

\section{Conceito de Branch-and-Bound}

\section{Funcionamento da Divisão do Espaço de Busca (Branching)}

\section{Conceito de Limites e Critérios de Poda (Bounding)}

\section{Aplicação da Estratégia em um Exemplo de Otimização}

\section{Apresentação de um Esquema ou Pseudocódigo Simplificado}

\section{Comparação entre Backtracking e Branch-and-Bound}

\section{Comparação Geral das Estratégias de Resolução e Busca (Parte I)}

% ----------------------------------------------------------
% PARTE II
% ----------------------------------------------------------
\part{Aproximação e Qualidade das Soluções}
% ----------------------------------------------------------

% ----------------------------------------------------------
% Capítulo 5 (Integrante 5)
% ----------------------------------------------------------
\chapter{Algoritmos Aproximados e Cobertura de Vértices}

\section{Conceito de Algoritmo Aproximado}

Quando lidamos com problemas de otimização pertencentes à classe NP-difícil, a impossibilidade prática de obter soluções exatas em tempo hábil motivou o desenvolvimento de uma classe rigorosa de métodos conhecida como \textbf{Algoritmos de Aproximação}.

\subsection{Definição de Algoritmo de Aproximação}

Algoritmos de aproximação são métodos computacionais eficientes --- isto é, que executam obrigatoriamente em \textbf{tempo polinomial} --- projetados para resolver problemas de otimização difíceis (geralmente NP-difíceis), nos quais encontrar a solução exata demandaria um tempo de processamento exponencial ou proibitivo. Diferentemente das heurísticas comuns, um algoritmo de aproximação possui uma prova matemática formal que atesta a qualidade da solução gerada para qualquer instância de entrada.

\subsection{Relação com a Busca pela Solução Ótima e Motivação}

A relação fundamenta-se em um compromisso (\textit{trade-off}) rigorosamente controlado: abre-se mão da certeza de alcançar o resultado matematicamente perfeito (a solução ótima global) em troca de velocidade de processamento (tempo polinomial), mas mantém-se o controle estrito sobre a margem de erro máxima por meio de uma \textbf{garantia matemática}. Assim, embora o algoritmo não entregue necessariamente o valor ótimo $OPT$, ele utiliza propriedades estruturais do próprio ótimo (como relaxações lineares ou limites inferiores) para assegurar que a resposta obtida nunca se afaste do ideal além de um fator conhecido.

O principal motivo para utilizar uma solução aproximada reside na barreira da complexidade computacional: achar a solução ótima em problemas NP-difíceis exige explorar espaços de busca que crescem de forma exponencial ou fatorial, tornando o processamento inviável para entradas reais de médio e grande porte. Os algoritmos de aproximação reduzem esse tempo drasticamente para uma escala polinomial, oferecendo a segurança adicional --- inexistente nas heurísticas puras --- de que o sistema jamais produzirá uma solução arbitrariamente ruim, mesmo nos piores cenários de entrada.

\section{Diferença entre Algoritmos Exatos, Heurísticos e Aproximados}

Na Análise de Algoritmos, a distinção fundamental entre os métodos exatos, heurísticos e aproximados reside no compromisso entre o tempo de execução e a garantia matemática de qualidade da solução:

\begin{itemize}
    \item \textbf{Algoritmos Exatos:} Exploram sistematicamente o espaço de busca até encontrar e comprovar matematicamente a solução ótima global ($C = C^*$). Contudo, para problemas NP-difíceis, exigem tempo exponencial ou fatorial no pior caso, tornando-se inviáveis para instâncias de grande escala.
    \item \textbf{Algoritmo que apenas procura uma solução boa (Heurística):} Baseia-se na ``intuição'', em regras práticas ou em tentativas empíricas inteligentes para achar uma solução rápida e satisfatória. Contudo, ele \textbf{não oferece garantias formais}: sob uma determinada condição favorável, pode encontrar a solução perfeita; em outra configuração de dados, pode entregar uma solução péssima sem que o usuário tenha meios de perceber o prejuízo.
    \begin{itemize}
        \item \textit{Exemplo ilustrativo:} É equivalente a pedir uma recomendação de restaurante para um amigo. Geralmente a refeição será boa e você economizará tempo de pesquisa, mas não existe nenhuma garantia científica ou contratual de que aquele é o melhor custo-benefício da cidade ou de que o prato não virá ruim naquele dia.
    \end{itemize}
    \item \textbf{Algoritmo com garantia matemática de qualidade (Algoritmo de Aproximação):} Também busca soluções rápidas em tempo polinomial, mas é construído sobre propriedades matemáticas estritas que provam o quão próxima sua resposta está da solução ótima. O desenvolvedor conhece de antemão o \textbf{pior cenário possível}. Por exemplo, um algoritmo $2$-aproximado para um problema de minimização garante matematicamente que o custo final jamais ultrapassará o dobro do custo da solução ótima ($C \le 2 \cdot C^*$), não importa o quão difícil seja a instância.
    \begin{itemize}
        \item \textit{Exemplo ilustrativo:} É equivalente a contratar um serviço de logística expressa com cláusula contratual auditada que garante que, mesmo no pior cenário possível (trânsito intenso ou chuva), sua encomenda será entregue em no máximo 48 horas (sendo o tempo ideal de um voo dedicado igual a 24 horas). O usuário tem uma segurança previsível e um teto máximo de atraso garantido.
    \end{itemize}
\end{itemize}

\section{Conceito de Razão de Aproximação e Interpretação da Qualidade}

Na Análise de Algoritmos, a \textbf{garantia de aproximação} (ou garantia formal de qualidade) é uma prova matemática rigorosa, estabelecida por meio de teoremas, que assegura um limite máximo para o erro de um algoritmo em qualquer cenário possível, independentemente do quão complexa ou atípica seja a entrada de dados.

Em sentido técnico e prático, dizer que uma solução é ``próxima da solução ideal'' significa que o algoritmo encontrou uma resposta viável e satisfatória em um tempo de execução aceitável, cujo custo $C$ está matematicamente limitado a uma pequena fração ou múltiplo do custo da solução ótima verdadeira $C^*$. 

Formalmente, diz-se que um algoritmo possui um \textbf{fator (ou razão) de aproximação} $\alpha$ se, para qualquer instância do problema, o valor da solução entregue difere do valor ótimo por, no máximo, um fator $\alpha$. Por exemplo, em um problema de minimização, uma $\alpha$-aproximação com $\alpha = 1,5$ garante que o custo obtido será, no pior caso, apenas $50\%$ maior que o ideal ($C \le 1,5 \cdot C^*$). Dessa forma, sacrifica-se parcialmente a precisão absoluta para obter o melhor resultado comprovadamente próximo do ótimo sem incorrer em tempos exponenciais.

\subsection{Demonstração por Meio de um Exemplo Simples}

\section{Definição do Problema de Cobertura de Vértices (Vertex Cover)}

\section{Aplicação de um Algoritmo de Aproximação em um Pequeno Grafo}

\section{Explicação da Garantia de Aproximação do Algoritmo Utilizado}

% ----------------------------------------------------------
% Capítulo 6 (Integrante 6)
% ----------------------------------------------------------
\chapter{Algoritmos Aproximados para o TSP}

\section{Dificuldades de Encontrar a Solução Ótima para o TSP}

\section{Diferença entre Encontrar uma Solução Viável e uma Solução Ótima}

Em uma operação logística real envolvendo o Problema do Caixeiro-Viajante, \textbf{uma rota viável e suficientemente boa (subótima) atende perfeitamente à necessidade da empresa}, não sendo necessário (e muitas vezes nem desejável) aguardar pela melhor rota absoluta possível (solução ótima). Essa diferença prática entre buscar apenas uma solução viável de boa qualidade e exigir a solução ótima justifica-se por três fatores fundamentais:

\begin{enumerate}
    \item \textbf{Restrição de tempo de resposta (\textit{Time-to-Decision}):} Uma transportadora precisa despachar seus veículos diariamente logo após o fechamento dos pedidos. É muito mais vantajoso utilizar um método rápido que leva 2 segundos para encontrar uma rota apenas $2\%$ ou $3\%$ mais longa do que a ideal, do que paralisar a frota por horas ou dias aguardando que o servidor calcule uma economia irrisória de poucos quilômetros.
    \item \textbf{Dinamismo e incerteza do ambiente real:} O mundo físico não é estático. A ``rota matematicamente perfeita'' calculada no escritório pode rapidamente se tornar a pior opção na prática devido a variáveis imprevisíveis, como acidentes de trânsito, congestionamentos repentinos, condições climáticas adversas, obras na pista ou atrasos no recebimento da mercadoria pelo cliente.
    \item \textbf{Retorno financeiro decrescente:} O custo computacional e de infraestrutura de TI (energia elétrica, servidores de alto desempenho e tempo de processamento) exigido para refinar uma solução de $98\%$ de eficiência para $100\%$ de eficiência frequentemente supera o valor financeiro do combustível que seria economizado naquele pequeno trecho otimizado.
\end{enumerate}

\section{Apresentação de um Algoritmo Aproximado para o TSP}

\section{Demonstração do Funcionamento do Algoritmo em um Exemplo Pequeno}

\section{Explicação da Garantia de Aproximação Oferecida pelo Método}

\section{Condições Necessárias para a Validade da Garantia (Desigualdade Triangular no TSP Métrico)}

% ----------------------------------------------------------
% Capítulo 7 (Integrante 7)
% ----------------------------------------------------------
\chapter{Set Cover, Bin Packing e Comparação entre Estratégias}

\section{Definição do Problema Set Cover}

\section{Definição do Problema Bin Packing}

\section{Exemplo de Aplicação de Estratégia Gulosa ou Aproximada}

Um caso clássico de aplicação prática derivado diretamente do Problema de Cobertura de Conjuntos (\textit{Set Cover}) --- que é comprovadamente NP-difícil --- é o \textbf{Problema da Cobertura Máxima Orçada} (\textit{Budgeted Maximum Coverage Problem}), como na instalação de antenas de internet para cobrir moradores respeitando um orçamento limitado.

\subsection{Funcionamento da Estratégia Gulosa de Custo-Benefício}

A ideia central da estratégia consiste em escolher iterativamente o local de antena (subconjunto) que oferece o \textbf{melhor custo-benefício marginal imediato}, contabilizando apenas os moradores que ainda não possuem sinal de internet. O funcionamento estrutura-se nos seguintes passos:

\begin{enumerate}
    \item \textbf{Inicialização:} Define-se o orçamento total disponível $B$, inicializa-se o conjunto de pessoas já atendidas como vazio ($C = \emptyset$) e mantém-se a lista de locais candidatos disponíveis para instalação.
    \item \textbf{Cálculo do Ganho Marginal:} A cada iteração, para cada local candidato $i$ ainda não selecionado, calcula-se quantas \textbf{pessoas novas} ($|S_i \setminus C|$) aquela antena passaria a atender, descontando-se os indivíduos que já estão cobertos por antenas instaladas em passos anteriores.
    \item \textbf{Avaliação da Razão de Eficiência:} Para cada local $i$, divide-se o número de novas pessoas cobertas pelo custo financeiro de instalação naquele ponto ($c_i$), obtendo a razão de custo-benefício:
    \begin{equation}
        R_i = \frac{|S_i \setminus C|}{c_i}
    \end{equation}
    \item \textbf{Seleção e Verificação Orçamentária:} Seleciona-se o local que apresenta a maior razão $R_i$ e verifica-se se o seu custo $c_i$ cabe no orçamento restante. Se couber, o local é adicionado à solução final, seu custo é descontado do orçamento, as novas pessoas são unidas ao conjunto de atendidos ($C$) e o local é removido dos disponíveis. Caso o local mais eficiente estoure o orçamento restante, ele é descartado e testa-se o próximo melhor candidato da lista.
    \item \textbf{Critério de Parada:} O laço de repetição continua até que o orçamento disponível se esgote ou não restem locais candidatos viáveis.
\end{enumerate}

\subsection{Vantagem de Velocidade e Garantia de Qualidade}

Enquanto o método de força bruta precisaria testar todos os subconjuntos possíveis de locais que respeitam o orçamento --- o que gera uma explosão combinatória da ordem de $O(2^m)$ para $m$ locais candidatos ---, a estratégia gulosa atua como um gerenciador incremental. Ela avalia apenas o estado atual, toma uma decisão definitiva por etapa e nunca retrocede (\textit{no backtracking}). Com isso, no pior caso, o algoritmo realiza no máximo $m$ iterações, inspecionando os locais restantes em cada passo, o que reduz a complexidade para um \textbf{tempo polinomial quadrático} $O(m^2)$.

\textbf{A estratégia não garante encontrar a solução ótima absoluta.} Como ela toma decisões localmente ótimas (escolhendo o melhor índice imediato naquele exato instante), uma antena escolhida no início pode consumir uma fatia do orçamento que impediria a compra combinada de duas outras antenas ligeiramente menos eficientes individualmente, mas que juntas cobririam um público maior. Entretanto, por se tratar de uma função submodular (retornos decrescentes de cobertura), essa estratégia gulosa modificada configura um \textbf{algoritmo de aproximação com desempenho comprovado}: a literatura demonstra que variações desse método garantem matematicamente atingir pelo menos $(1 - 1/e) \approx 63,2\%$ da cobertura da solução ótima global.

Para mensurar a eficiência prática e teórica da solução encontrada pelo sistema, são fundamentais as seguintes métricas e informações:

\begin{itemize}
    \item \textbf{Índice de Custo Efetivo por Usuário Atendido:} A razão consolidada entre o total de pessoas beneficiadas e o montante financeiro efetivamente investido.
    \item \textbf{Taxa de Utilização do Orçamento e Cobertura Demográfica:} O percentual da população total da região que passou a ter acesso à internet e a verificação de eventual saldo ocioso no orçamento.
    \item \textbf{Comparativo com um Limite Superior Matemático (\textit{Upper Bound}):} O resultado da relaxação do problema (por exemplo, permitindo antenas ``fracionárias'' via Programação Linear) ou a cota teórica de aproximação, permitindo saber exatamente o quão perto do ideal teórico (desconhecido) o algoritmo chegou.
\end{itemize}

\subsection{Exemplo Numérico Passo a Passo}

\section{Comparação entre Heurísticas e Algoritmos Aproximados}

A \autoref{tab:comparativo_heuristicas_aproximacao} sintetiza as diferenças conceituais, teóricas e práticas entre as abordagens heurísticas e os algoritmos de aproximação:

\begin{table}[htb]
\centering
\caption{\label{tab:comparativo_heuristicas_aproximacao}Comparativo técnico entre Heurísticas e Algoritmos de Aproximação.}
\small
\renewcommand{\arraystretch}{1.4}
\begin{tabular}{|>{\raggedright\arraybackslash}p{3.0cm}|>{\raggedright\arraybackslash}p{5.5cm}|>{\raggedright\arraybackslash}p{5.5cm}|}
\hline
\rowcolor{gray!20}
\textbf{Critério} & \textbf{Heurística} & \textbf{Algoritmo de Aproximação} \\ \hline
\textbf{Objetivo Principal} & Encontrar uma solução rápida e satisfatória (``boa o suficiente'') para problemas computacionais muito complexos ou de grande escala. & Encontrar soluções de boa qualidade (comprovadamente próximas da ótima) em tempo polinomial para problemas de otimização difíceis (NP-difíceis). \\ \hline
\textbf{Garantem a solução ótima?} & Não. & Não. \\ \hline
\textbf{Possuem garantia formal de qualidade?} & Não. O erro no pior caso pode ser arbitrariamente grande. & Sim. Possuem um fator ou razão matemática de aproximação comprovada por teorema. \\ \hline
\textbf{Principal vantagem} & Redução drástica da complexidade computacional e do tempo de resposta, aliada à facilidade de adaptação a regras de negócio complexas. & Segurança e previsibilidade: garantia matemática da qualidade da solução dentro de um tempo de execução rápido (polinomial). \\ \hline
\textbf{Principal limitação} & Falta de garantia de otimalidade e de limite de erro (risco de ficar preso em ótimos locais ruins sem detecção). & Rigidez matemática: nem todo problema complexo admite uma aproximação eficiente devido a barreiras teóricas de \textit{inaproximabilidade} (como o Teorema PCP). \\ \hline
\textbf{Exemplo de aplicação} & Algoritmos de busca informada em aplicativos de GPS (como Waze e Google Maps) e heurística do Vizinho Mais Próximo. & Algoritmo de Christofides ($1,5$-aproximação) para planejamento de rotas no Caixeiro-Viajante Métrico e Cobertura de Vértices. \\ \hline
\end{tabular}
\end{table}

\section{Discussão sobre Qualidade das Soluções, Tempo de Execução e Garantias Teóricas}

A adoção de métodos heurísticos ou aproximados no projeto de algoritmos envolve um balanço crítico entre o custo computacional de execução e a precisão matemática do resultado obtido. No caso das heurísticas, destacam-se duas vantagens fundamentais e duas limitações intrínsecas:

\subsection{Vantagens do Uso de Heurísticas}

\begin{enumerate}
    \item \textbf{Viabilidade Temporal e Escalabilidade para Grandes Instâncias:} A principal vantagem de uma heurística é a sua capacidade de fornecer respostas válidas em tempo polinomial baixo (frequentemente linear $O(n)$, log-linear $O(n \log n)$ ou quadrático $O(n^2)$) para problemas cuja solução exata exigiria tempo exponencial $O(2^n)$ ou fatorial $O(n!)$. Na prática, isso permite que sistemas computacionais escalem para lidar com bancos de dados reais contendo dezenas de milhares de variáveis --- como o roteamento logístico de uma metrópole ou a alocação de processos em servidores em nuvem --- entregando decisões em segundos ou milissegundos, onde algoritmos exatos travariam por exaustão de tempo ou de memória RAM.
    \item \textbf{Simplicidade de Modelagem e Flexibilidade de Adaptação:} Algoritmos exatos (como Programação Linear Inteira ou \textit{Branch-and-Bound}) costumam exigir formulações matemáticas rígidas e provas formais complexas que se desfazem facilmente quando uma nova regra de negócio é adicionada ao problema. As heurísticas, por outro lado, são baseadas em regras intuitivas e construtivas. Isso torna sua implementação em código mais direta e permite que o desenvolvedor incorpore restrições complexas do mundo real (como janelas de horário dinâmicas, penalidades não lineares ou preferências do usuário) apenas ajustando a função de avaliação ou os critérios de escolha local.
\end{enumerate}

\subsection{Limitações do Uso de Heurísticas}

\begin{enumerate}
    \item \textbf{Ausência de Garantia de Otimalidade e de Limite de Erro (Estagnação em Ótimos Locais):} Como as heurísticas descartam grandes porções do espaço de busca e tomam decisões baseadas em informações restritas (visão local ou míope), elas não garantem encontrar a solução ótima global. Mais grave ainda do ponto de vista teórico: uma heurística pura não oferece uma cota matemática formal sobre o quão distante a resposta entregue está da solução ideal. O algoritmo pode ficar preso em uma ``armadilha'' topológica (um ótimo local) e retornar uma solução que é $5\%$ pior que a ótima em um dia, mas $80\%$ pior no dia seguinte, sem que o sistema consiga detectar essa degradação por conta própria.
    \item \textbf{Alta Sensibilidade aos Dados de Entrada e Especificidade de Domínio:} O desempenho qualitativo de uma heurística específica costuma ser fortemente dependente da estrutura ou distribuição dos dados de entrada (bem como da escolha de parâmetros iniciais, como o vértice de partida). Uma regra prática projetada sob medida para funcionar muito bem no Problema do Caixeiro-Viajante com distâncias euclidianas planas pode produzir resultados desastrosos se aplicada a uma rede com custos assimétricos (como tarifas aéreas) ou a um problema de escalonamento de tarefas. Assim, heurísticas clássicas carecem de generalidade, exigindo reengenharia e baterias extensas de testes empíricos sempre que o cenário do problema sofre alterações estruturais.
\end{enumerate}

\section{Identificação de Situações Mais Adequadas para Cada Estratégia no Mundo Real}

A seguir, são analisadas três situações reais da indústria e da computação nas quais o uso de heurísticas e algoritmos de aproximação é indispensável para viabilizar a operação.

\subsection{Situação 1: Organização Industrial e Escalonamento de Tarefas (\textit{Job Shop Scheduling})}

\begin{itemize}
    \item \textbf{Problema a ser resolvido:} Determinar a ordem sequencial e o momento exato em que um conjunto de tarefas (ou lotes de produtos) deve ser processado em um conjunto de máquinas industriais (ou núcleos de um processador), respeitando restrições técnicas de precedência e prazos, de modo a minimizar o tempo total de conclusão da produção (\textit{makespan}).
    \item \textbf{Por que encontrar a solução ideal é difícil:} Devido ao crescimento exponencial e fatorial das combinações possíveis. À medida que o número de tarefas ($n$) e de máquinas ($m$) aumenta, a quantidade de escalonamentos distintos cresce na ordem de $(n!)^m$. Para dezenas de tarefas e máquinas, testar todas as sequências torna-se matematicamente intratável.
    \item \textbf{Como uma heurística ou algoritmo de aproximação ajuda:} Transformam um problema intratável em um processo resolvido em segundos. Algoritmos de aproximação clássicos (como o \textit{Escalonamento em Lista de Graham}, que possui garantia de aproximação de $2 - \frac{1}{m}$) ou heurísticas de despacho como \textit{Menor Tempo de Processamento Primeiro} (SPT) alocam as tarefas imediatamente assim que uma máquina fica livre, entregando um cronograma altamente produtivo e próximo do ideal sem paralisar a fábrica.
\end{itemize}

\subsection{Situação 2: Roteamento de Veículos e Logística de Transporte (\textit{Vehicle Routing Problem -- VRP})}

\begin{itemize}
    \item \textbf{Problema a ser resolvido:} Encontrar o conjunto de rotas mais eficiente para uma frota de múltiplos veículos atender a centenas de clientes espalhados por diferentes regiões geográficas, utilizando o mínimo possível de combustível e tempo, sem violar a capacidade de carga de cada caminhão.
    \item \textbf{Por que encontrar a solução ideal é difícil:} Além de generalizar o Problema do Caixeiro-Viajante (herdando sua explosão combinatória NP-difícil), o problema incorpora múltiplas restrições simultâneas do mundo real --- como limites de peso/volume da frota, janelas de horário de recebimento dos clientes e variações dinâmicas de trânsito --- que inviabilizam algoritmos exatos em escala comercial.
    \item \textbf{Como uma heurística ou algoritmo de aproximação ajuda:} Atuam diretamente nas duas maiores exigências da operação: velocidade de tomada de decisão e respeito às restrições rígidas. Heurísticas construtivas consagradas, como o \textbf{Algoritmo das Economias de Clarke e Wright}, combinadas com meta-heurísticas de busca local (como \textit{2-opt} ou \textit{Simulated Annealing}), reduzem um cálculo que levaria séculos para poucos minutos, gerando rotas logísticas enxutas e prontas para execução diária.
\end{itemize}

\subsection{Situação 3: Armazenamento e Computação em Nuvem (\textit{Bin Packing} e Problema da Mochila)}

\begin{itemize}
    \item \textbf{Problema a ser resolvido:} Distribuir e alocar uma quantidade massiva de máquinas virtuais, contêineres de aplicações ou arquivos digitais dentro do menor número possível de servidores físicos em um \textit{Data Center}, otimizando o consumo de energia elétrica e o uso de hardware.
    \item \textbf{Por que encontrar a solução ideal é difícil:} Trata-se do clássico problema de empacotamento (\textit{Bin Packing}) em múltiplas dimensões. Diferentemente de caixas físicas onde avalia-se apenas o volume ou peso, na computação em nuvem cada tarefa consome simultaneamente múltiplas dimensões de recursos (CPU, memória RAM, largura de banda de rede e leitura de disco), além de sofrer alterações dinâmicas a cada segundo com novas requisições de usuários.
    \item \textbf{Como uma heurística ou algoritmo de aproximação ajuda:} O problema é solucionado em tempo real aplicando regras lógicas rápidas de empacotamento, sem necessidade de testar todas as permutações de servidores. Algoritmos de aproximação gulosos como o \textbf{\textit{First-Fit Decreasing} (FFD)} --- que ordena os processos do maior para o menor e aloca cada um no primeiro servidor com espaço disponível --- executam em tempo $O(n \log n)$ e garantem matematicamente utilizar um número de servidores muito próximo do mínimo teórico possível.
\end{itemize}

% ----------------------------------------------------------
% Finaliza a parte no bookmark do PDF
% ----------------------------------------------------------
\phantompart

% ----------------------------------------------------------
% Considerações Finais
% ----------------------------------------------------------
\chapter*[Considerações Finais]{Considerações Finais}
\addcontentsline{toc}{chapter}{Considerações Finais}
% ----------------------------------------------------------

% ----------------------------------------------------------
% ELEMENTOS PÓS-TEXTUAIS
% ----------------------------------------------------------
\postextual
% ----------------------------------------------------------

% ----------------------------------------------------------
% Referências bibliográficas
% ----------------------------------------------------------
\chapter*[Referências]{Referências}
\addcontentsline{toc}{chapter}{Referências}

\noindent
BARBOSA, Marco Aurélio Lima. \textbf{Ciclos Hamiltonianos e o Problema do Caixeiro Viajante (Notas de Aula de Grafos)}. Departamento de Informática -- Universidade Estadual de Maringá (DIN-UEM), 2022. Disponível em: \url{https://malbarbo.pro.br/arquivos/2022/6898/12-ciclos-hamiltonianos-e-o-problema-do-caixeiro-viajante.pdf}.

\noindent
CORMEN, T. H. et al. \textbf{Algoritmos: Teoria e Prática}. 3. ed. Rio de Janeiro: Elsevier, 2012. (Capítulo 34: NP-Completude e Seção 35.2: O Problema do Caixeiro-Viajante).

\noindent
FEOFILOFF, Paulo. \textbf{Algoritmos Gulosos}. Projeto de Algoritmos em C -- Instituto de Matemática e Estatística da Universidade de São Paulo (IME-USP). Disponível em: \url{https://www.ime.usp.br/~pf/analise_de_algoritmos/aulas/guloso.html}.

\noindent
FEOFILOFF, Paulo. \textbf{Minicurso de Análise de Algoritmos}. Departamento de Ciência da Computação -- Instituto de Matemática e Estatística da Universidade de São Paulo (IME-USP), 2009. (Seção 1.2: Consumo de tempo de algoritmos e Comparação assintótica de funções). Disponível em: \url{https://www.ime.usp.br/~pf/livrinho-AA/downloads/AA-BOOKLET.pdf}.

\noindent
FEOFILOFF, Paulo. \textbf{O Problema do Caixeiro Viajante}. Projeto de Algoritmos em C -- Instituto de Matemática e Estatística da Universidade de São Paulo (IME-USP). Disponível em: \url{https://www.ime.usp.br/~pf/analise_de_algoritmos/aulas/caixeiro.html}.

\noindent
FEOFILOFF, Paulo; KOHAYAKAWA, Yoshiharu; WAKABAYASHI, Yoshiko. \textbf{Uma Introdução Sucinta à Teoria dos Grafos}. Instituto de Matemática e Estatística da Universidade de São Paulo (IME-USP), 2011. Disponível em: \url{https://www.ime.usp.br/~pf/teoriadosgrafos/texto/TeoriaDosGrafos.pdf}.

\noindent
FERNANDES, Cristina Gomes. \textbf{MAC5758 / MAC0466 -- Algoritmos de Aproximação}. Instituto de Matemática e Estatística da Universidade de São Paulo (IME-USP), 2021. Disponível em: \url{https://www.ime.usp.br/~cris/aprox/}.

\noindent
GOLDBARG, Marco Cesar; LUNA, Henrique Pacca Loureiro. \textbf{Otimização Combinatória e Programação Linear: Modelos e Algoritmos}. 2. ed. Rio de Janeiro: Elsevier, 2005. (Capítulo sobre Heurísticas e Meta-heurísticas).

\noindent
MIYAZAWA, Flávio Keidi. \textbf{Introdução à Otimização Combinatória}. Instituto de Computação da Universidade Estadual de Campinas (IC-Unicamp). Disponível em: \url{https://www.ic.unicamp.br/~fkm/lectures/intro-otimizacao.pdf}.

\noindent
MIYAZAWA, Flávio Keidi. \textbf{Projetos em Algoritmos de Aproximação}. Instituto de Computação da Universidade Estadual de Campinas (IC-Unicamp). Disponível em: \url{https://www.ic.unicamp.br/~fkm/problems/algaprox.html}.

\noindent
PARPINELLE, Rafael Stubs. \textbf{Configuração Automática de Algoritmos Heurísticos de Otimização}. Centro de Educação Superior do Alto Vale do Itajaí -- Universidade do Estado de Santa Catarina (UDESC). Disponível em: \url{https://www.udesc.br/arquivos/ceavi/id_cpmenu/5818/Configura__o_autom_tica_de_algoritmos_heur_sticos_de_otimiza__o_17315186624122_5818.pdf}.

\noindent
RUSSELL, Stuart; NORVIG, Peter. \textbf{Inteligência Artificial}. 3. ed. Rio de Janeiro: Elsevier, 2013. (Capítulo 3, Seção 3.5: Estratégias de Busca Informada / Heurística e Capítulo 4: Busca Local em Espaços de Estados).

\noindent
SCIENCEDIRECT TOPICS. \textbf{Heuristic Algorithm -- Computer Science}. Elsevier. Disponível em: \url{https://www.sciencedirect.com/topics/computer-science/heuristic-algorithm}.

\noindent
SOUZA, Marcone Jamilson Freitas. \textbf{Inteligência Computacional para Otimização (Notas de Aula -- Técnicas Meta-heurísticas para Otimização Combinatória)}. Departamento de Computação -- Universidade Federal de Ouro Preto (DECOM-UFOP). Disponível em: \url{http://www.decom.ufop.br/marcone/Disciplinas/InteligenciaComputacional/InteligenciaComputacional.htm}.

\noindent
TECHVIDVAN. \textbf{Greedy Algorithm -- Introduction, Methods and Examples}. Disponível em: \url{https://techvidvan.com/tutorials/greedy-algorithm/}.

\noindent
WIKIPEDIA. \textbf{Maximum Coverage Problem}. Wikimedia Foundation. Disponível em: \url{https://en.wikipedia.org/wiki/Maximum_coverage_problem}.

\noindent
WIKIPÉDIA. \textbf{Algoritmo de Aproximação}. Wikimedia Foundation. Disponível em: \url{https://pt.wikipedia.org/wiki/Algoritmo_de_aproxima%C3%A7%C3%A3o}.

\noindent
ZIVIANI, Nivio. \textbf{Projeto de Algoritmos com Implementações em Pascal e C}. 3. ed. São Paulo: Cengage Learning, 2011. (Capítulo sobre Algoritmos Aproximativos e Problemas NP-Completos).

\end{document}
